%% ============================================================================
%%  中文简历模板（skeleton）— ai-job-search DeepSeek 版附带
%% ============================================================================
%%  用途：中文岗位投递用的 1 页简历骨架。上游自带的 moderncv / cover.cls 不含
%%        中文字形，直接填中文会缺字，所以中文投递走这个模板。
%%
%%  编译：xelatex --interaction=nonstopmode --halt-on-error template.tex
%%        （不要用 lualatex 走 moderncv 那条路；本模板用 ctexart + xeCJK）
%%
%%  字体：默认用 Windows 自带的「微软雅黑」。Linux/macOS 请改 \setCJKmainfont：
%%        - macOS:  PingFang SC
%%        - Linux:  Noto Sans CJK SC / Source Han Sans SC
%%        - 或者统一用 TeX Live 自带的 FandolSong / FandolHei
%%
%%  占位符：所有 【...】 都要替换。不要留着占位符就去投。
%%  英文投递请改回上游的 cv/main_example.tex（moderncv, lualatex）。
%%
%%  页数：应届生 / 经历较少 → 1 页；有多年履历 → 2 页。宁可删，不要挤。
%% ============================================================================

\documentclass[10.5pt,a4paper]{ctexart}

\usepackage[scale=0.86]{geometry}
\usepackage{xeCJK}
\usepackage{fontspec}
\usepackage{titlesec}
\usepackage{enumitem}
\usepackage[hidelinks]{hyperref}
\usepackage{xcolor}

\definecolor{accent}{HTML}{1F4E79}

% ---- 字体（按你的系统改这一行即可）----
\setCJKmainfont{Microsoft YaHei}[BoldFont=Microsoft YaHei Bold]
\setCJKsansfont{Microsoft YaHei}
\setmainfont{Times New Roman}

\titleformat{\section}{\large\bfseries\color{accent}}{}{0em}{}[\vspace{-2pt}\rule{\linewidth}{0.6pt}]
\titlespacing*{\section}{0pt}{9pt}{5pt}
\setlist[itemize]{leftmargin=1.2em,itemsep=1pt,topsep=2pt,parsep=0pt}
\pagestyle{empty}

%% 一条经历：{\entry{单位/学校}{时间}{职位/专业}{地点}}
\newcommand{\entry}[4]{%
  \noindent\textbf{#1}\hfill #2\\[1pt]
  {\small #3 \hfill #4}\par\vspace{3pt}}

\begin{document}

% ============================== 抬头 ==============================
\begin{center}
  {\Huge\bfseries 【姓名】}\\[4pt]
  {\small 求职意向：【目标岗位】\quad|\quad【性别】·【民族】·【政治面貌】·【出生年月】}\\[2pt]
  {\small 电话 【手机号】\quad|\quad 邮箱 【邮箱】\quad|\quad【所在城市】}
\end{center}

\vspace{2pt}

% ============================ 教育背景 ============================
\section{教育背景}
\entry{【学校】}{【入学年月】 -- 【毕业年月】}{【专业】（【学历】）}{【城市】}
主修课程：【课程 1】、【课程 2】、【课程 3】、【课程 4】。【一句话补充：成绩排名 / 奖学金 / 在校实践。】

% ========================== 实习与工作经历 ==========================
% 应届生：实习在前、校内职务在后（或合并）。有多段实习就复制一组。
\section{实习与工作经历}
\entry{【公司名】}{【开始年月】 -- 【结束年月】}{【岗位】}{【城市】}
\begin{itemize}
  \item 【职责/成果 1：动词开头，写清做了什么、对象是谁、结果如何】
  \item 【职责/成果 2】
  \item 【职责/成果 3】
  \item 【职责/成果 4】
\end{itemize}

\entry{【组织/公司名】}{【开始年月】 -- 【结束年月】}{【职务】}{【城市】}
\begin{itemize}
  \item 【职责/成果 1】
  \item 【职责/成果 2】
  \item 【职责/成果 3】
\end{itemize}

% ============================ 技能与证书 ============================
\section{技能与证书}
\begin{itemize}
  \item \textbf{办公软件}：【如 Word、Excel、PowerPoint；具体到能做什么】
  \item \textbf{语言能力}：【中文（母语）；英语 CET-4 / CET-6，读写/口语水平】
  \item \textbf{综合能力}：【如沟通协调、活动组织、文字写作、档案管理】
\end{itemize}

% ============================= 自我评价 =============================
% 可选。没有独特内容可以整段删掉，不要写空话套话。
\section{自我评价}
【2-3 句，说明你的工作方式与能给团队带来什么。不要写"吃苦耐劳""团队精神"这类无信息量的词。】

\end{document}
